\documentclass[conference]{IEEEtran}

% =========================================================
% PLANTILLA LIGERA PARA OVERLEAF FREE
% Compilador recomendado: pdfLaTeX
% =========================================================

\usepackage[utf8]{inputenc}
\usepackage[T1]{fontenc}
\usepackage[spanish,es-nodecimaldot]{babel}
\usepackage{lmodern}
\usepackage{microtype}

\usepackage{amsmath,amssymb}
\usepackage{graphicx}
\usepackage{booktabs}
\usepackage{multirow}
\usepackage{array}
\usepackage{url}
\usepackage{cite}
\usepackage{xcolor}
\usepackage{listings}
\usepackage{algorithm}
\usepackage{algpseudocode}
\usepackage{hyperref}

\definecolor{CodeBackground}{HTML}{F4F6F7}
\definecolor{CodeKeyword}{HTML}{005F73}
\definecolor{CodeComment}{HTML}{6C757D}
\definecolor{CodeString}{HTML}{BB3E03}

\hypersetup{
    colorlinks=true,
    linkcolor=black,
    citecolor=black,
    urlcolor=blue
}

\lstdefinestyle{paperstyle}{
    backgroundcolor=\color{CodeBackground},
    basicstyle=\ttfamily\footnotesize,
    keywordstyle=\color{CodeKeyword}\bfseries,
    commentstyle=\color{CodeComment}\itshape,
    stringstyle=\color{CodeString},
    numbers=left,
    numberstyle=\tiny,
    numbersep=5pt,
    frame=single,
    breaklines=true,
    showstringspaces=false,
    tabsize=4
}

\lstset{
    style=paperstyle,
    language=Python
}

% =========================================================
% DATOS DEL ARTÍCULO — EDITAR AQUÍ
% =========================================================

\title{
Título del Proyecto o Artículo Científico:
Subtítulo Descriptivo del Estudio
}

\author{
\IEEEauthorblockN{
Autor 1,
Autor 2,
Autor 3,
Autor 4
}
\IEEEauthorblockA{
\textit{Universidad Nacional de Colombia}\\
\textit{Sede Bogotá}\\
Bogotá D.C., Colombia\\
\{correo1, correo2, correo3, correo4\}@unal.edu.co
}
}
\begin{document}

\maketitle

% =========================================================
% RESUMEN
% =========================================================

\begin{abstract}
Este documento presenta una plantilla para la elaboración de proyectos,
artículos científicos y ponencias con estructura de congreso. El resumen debe
describir brevemente el problema, el objetivo, la metodología, los principales
resultados y la conclusión más importante. Se recomienda una extensión entre
150 y 250 palabras, dependiendo de las normas del evento.
\end{abstract}

\begin{IEEEkeywords}
inteligencia artificial, estructuras de datos, aprendizaje automático,
ingeniería de software, metodología experimental
\end{IEEEkeywords}

% =========================================================
% CONTENIDO
% =========================================================

\input{sections/01-introduccion}
\input{sections/02-trabajo-relacionado}
\input{sections/03-metodologia}
\input{sections/04-resultados}
\input{sections/05-discusion}
\input{sections/06-conclusiones}

% =========================================================
% AGRADECIMIENTOS
% =========================================================

\section*{Agradecimientos}

Los autores agradecen a la Universidad Nacional de Colombia, Sede Amazonia,
y a las personas o instituciones que apoyaron el desarrollo de este proyecto.

% =========================================================
% REFERENCIAS
% Sin BibTeX ni Biber para reducir tiempo de compilación.
% =========================================================

\begin{thebibliography}{00}

\bibitem{cormen}
T. H. Cormen, C. E. Leiserson, R. L. Rivest y C. Stein,
\textit{Introduction to Algorithms}, 4th ed.
Cambridge, MA, USA: MIT Press, 2022.

\bibitem{goodrich}
M. T. Goodrich, R. Tamassia y M. H. Goldwasser,
\textit{Data Structures and Algorithms in Python}.
Hoboken, NJ, USA: Wiley, 2013.

\bibitem{python}
Python Software Foundation,
``Python Documentation.''
[En línea]. Disponible en: \url{https://docs.python.org/3/}

\end{thebibliography}

\end{document}

